\documentclass[11pt,a4paper]{article}

\usepackage[utf8]{inputenc}
\usepackage[T1]{fontenc}
\usepackage{lmodern}
\usepackage[margin=1in]{geometry}
\usepackage{amsmath,amssymb,amsthm}
\usepackage{booktabs}
\usepackage{microtype}
\usepackage{xcolor}
\usepackage{enumitem}
\usepackage{algorithm}
\usepackage{algpseudocode}
\usepackage[hidelinks]{hyperref}
\usepackage{url}
\usepackage{listings}
\usepackage{caption}
\usepackage{multirow}

\definecolor{codebg}{rgb}{0.97,0.97,0.97}
\lstset{
  basicstyle=\ttfamily\small,
  backgroundcolor=\color{codebg},
  frame=single,
  framerule=0pt,
  showstringspaces=false,
  breaklines=true,
  language=C,
  keywordstyle=\color{blue!60!black},
  commentstyle=\color{green!50!black}\itshape,
  stringstyle=\color{red!60!black},
}

\newcommand{\Z}{\mathbb{Z}}
\newcommand{\N}{\mathbb{N}}
\newcommand{\Zq}{\Z_q}
\newcommand{\Rq}{R_q}
\newcommand{\NTT}{\mathrm{NTT}}
\newcommand{\iNTT}{\mathrm{NTT}^{-1}}
\newcommand{\Power}{\mathrm{Power2Round}}
\newcommand{\HighBits}{\mathrm{HighBits}}
\newcommand{\LowBits}{\mathrm{LowBits}}
\newcommand{\UseHint}{\mathrm{UseHint}}
\newcommand{\MakeHint}{\mathrm{MakeHint}}
\newcommand{\ExpandA}{\mathrm{ExpandA}}
\newcommand{\SHAKE}{\mathrm{SHAKE}}
\newcommand{\bigO}{\mathcal{O}}

\title{\Large\bfseries Specializing ML-DSA-44 Verification for Secure Boot on Cortex-M4 \\
       \large A Didactic Walk-through of Ten Progressive Optimizations}
\author{Oscar Reparaz}
\date{2026-05-11}

\begin{document}
\maketitle

\begin{abstract}
We present a custom verifier for ML-DSA-44 (NIST FIPS 204) targeted at
secure-boot deployments on the ARM Cortex-M4. Starting from the fastest
publicly available implementation (\texttt{pqm4}'s \texttt{m4f} backend),
we apply ten layered optimizations and reduce per-verify cost from
$1{,}261{,}555$ to $241{,}756$ instructions retired on the
Cortex-M4 ISA --- a $\mathbf{5.22\times}$ speedup. The optimizations exploit
two observations that hold for secure boot but not for general ML-DSA
verification:
(i)~the public key is fixed at provisioning time, so any quantity that
depends only on $\textsf{pk}$ may be precomputed and stored in flash;
(ii)~all verifier inputs are public, so the constant-time discipline of
\texttt{pqm4} is no longer load-bearing. The end-to-end verifier is
spec-conforming (still implements FIPS 204 Algorithm~8 byte-for-byte) and
verifies the same test vector mldsa-native produces. We give both
high-level intuition and the precise transformations applied at each step,
with measurements at every checkpoint reproducible from a small QEMU
harness.\smallskip

\noindent\textbf{Headline:} ML-DSA-44 signature verification on Cortex-M4
becomes $\sim$34$\times$ cheaper than pure-C P-256 ECDSA verify and
$\sim$51$\times$ cheaper than the hand-asm \texttt{p256-m}.
\end{abstract}

\section{Introduction}
\label{sec:intro}

The migration from elliptic-curve cryptography (ECC) to post-quantum
cryptography is now operational policy: NIST published the first three
post-quantum standards in August~2024, including FIPS~204 (ML-DSA, a
digital signature algorithm based on the Dilithium~\cite{dilithium} module
lattice construction). Secure-boot is among the first deployment domains
where the new standard matters: the public verification key is provisioned
into a device at manufacture, then the bootloader verifies a firmware
signature on every boot until end-of-life. Verifier cost is paid every
boot; key size is paid once.

ML-DSA's verify path is dominated by polynomial arithmetic over the ring
$\Rq = \Zq[X]/(X^{256}+1)$, $q = 2^{23} - 2^{13} + 1 = 8{,}380{,}417$, with
heavy use of the Number-Theoretic Transform (NTT). Compared to a
P-256 ECDSA verify on the same Cortex-M4, an off-the-shelf ML-DSA-44
verify in C is already $\sim$3--4$\times$ cheaper; with the hand-tuned
\texttt{pqm4}~\cite{pqm4} M4 asm it becomes $\sim$6$\times$ cheaper. This
paper asks how much further the cost can fall when one exploits the
\emph{fixed-public-key, all-public-data} setting that defines secure boot.

The contribution is twofold. First, we measure ten ML-DSA-44 verifier
variants on a single bare-metal Cortex-M4 testbench, with reproducible
instruction-retired counts. Second, and primarily, we deliver a didactic
exposition: each optimization is presented with its high-level
justification, the precise FIPS~204 step it touches, and its measured
impact, so a practitioner can decide which transformations apply to their
threat model and which do not. We are explicit about what we do
\emph{not} skip --- the FIPS-204-mandated checks on $z$'s norm and the
hint encoding remain in place --- because for secure boot the
correctness-of-verify property is the entire load-bearing assumption.

\section{Background}
\label{sec:background}

This section gives just enough lattice-cryptography to follow the
verification anatomy in Section~\ref{sec:verify-anatomy}. Readers fluent
in module-LWE and the Dilithium structure may skim.

\subsection{Module-LWE and the Ring \texorpdfstring{$\Rq$}{Rq}}

Let $q$ be a prime and $n$ a power of two. Define the
\emph{negacyclic} ring
\[
  \Rq \;:=\; \Zq[X] / (X^n + 1).
\]
Elements of $\Rq$ are polynomials of degree $<n$ with coefficients in
$\Zq$; multiplication is polynomial multiplication modulo $X^n + 1$, which
under our parameters ($n=256$, $q = 2^{23} - 2^{13} + 1$) admits a fast
$n\log n$ Number-Theoretic Transform.

The \emph{module} variant of the Learning-With-Errors problem (M-LWE)
fixes ranks $k,\ell \in \N$ and asks: distinguish
\[
  (\mathbf{A},\,\mathbf{A}\mathbf{s} + \mathbf{e})
  \quad\text{from}\quad
  (\mathbf{A},\,\mathbf{u})
\]
where $\mathbf{A} \in \Rq^{k\times\ell}$ is uniform, $\mathbf{s} \in
\Rq^{\ell}$ and $\mathbf{e} \in \Rq^{k}$ are short (coefficients in
$\{-\eta,\dots,\eta\}$ for a small $\eta$), and $\mathbf{u}$ is uniform in
$\Rq^{k}$. Both Kyber (now ML-KEM) and Dilithium (now ML-DSA) reduce their
security to M-LWE and its sibling M-SIS.

\subsection{Kyber (ML-KEM) at a Glance}

ML-KEM is a key-encapsulation mechanism. The public key is a M-LWE
sample $\mathbf{t} = \mathbf{A}\mathbf{s} + \mathbf{e}$; encapsulation
chooses small $\mathbf{r}, \mathbf{e}', e''$ and sends
$(\mathbf{u},v) = (\mathbf{A}^\top\mathbf{r} + \mathbf{e}',\;
\mathbf{t}^\top\mathbf{r} + e'' + \lceil q/2 \rceil \cdot m)$;
decapsulation recovers $m$ from $v - \mathbf{s}^\top \mathbf{u}$ since
the cross-term $\mathbf{e}^\top\mathbf{r} - \mathbf{s}^\top\mathbf{e}'
+ e''$ is small enough to fall on the correct side of the rounding
threshold. We do not optimize Kyber here, but the
$\mathbf{A} \mathbf{s} + \mathbf{e}$ pattern recurs.

\subsection{Dilithium (ML-DSA) Construction}
\label{sec:dilithium}

ML-DSA is a \emph{Fiat--Shamir-with-aborts} signature: a $\Sigma$-protocol
proving knowledge of a short M-LWE secret, made non-interactive by
deriving the challenge via a random oracle ($\SHAKE$-256). The
construction's three signature levels are parametrized as in
Table~\ref{tab:mldsa-params}.

\begin{table}[h]
\centering
\caption{ML-DSA parameter sets (FIPS 204, \S4 Table~1).}
\label{tab:mldsa-params}
\begin{tabular}{l c c c c c c c c c c}
\toprule
            & $q$       & $n$  & $d$ & $\tau$ & $\lambda$ & $\gamma_1$ & $\gamma_2$       & $(k,\ell)$ & $\eta$ & $\omega$ \\
\midrule
ML-DSA-44   & $2^{23}{-}2^{13}{+}1$ & 256 & 13 & 39 & 128 & $2^{17}$  & $(q{-}1)/88$ & $(4,4)$ & 2 & 80 \\
ML-DSA-65   & ,, & 256 & 13 & 49 & 192 & $2^{19}$  & $(q{-}1)/32$ & $(6,5)$ & 4 & 55 \\
ML-DSA-87   & ,, & 256 & 13 & 60 & 256 & $2^{19}$  & $(q{-}1)/32$ & $(8,7)$ & 2 & 75 \\
\bottomrule
\end{tabular}
\end{table}

For ML-DSA-44 a public key occupies $32 + k \cdot 320 = 1{,}312$~bytes
(rho plus the high-order $t_1$ of $\mathbf{t} = \mathbf{A}\mathbf{s}_1 +
\mathbf{s}_2$ truncated by $\Power$) and a signature occupies
$32 + \ell \cdot 576 + (\omega + k) = 2{,}420$~bytes (the challenge hash
$\tilde c$, the masked response $\mathbf{z}$, and the hint $\mathbf{h}$).

\subsection{The Number-Theoretic Transform}

Because $q-1$ is divisible by $2n = 512$, $\Zq$ contains a primitive
$2n$-th root of unity $\zeta$. The negacyclic NTT evaluates a polynomial
$a(X) \in \Rq$ at the $n$ odd powers of $\zeta$:
\[
  \widehat{a}_i \;=\; a(\zeta^{2i+1}) \pmod{q}, \quad i = 0,\dots,n-1.
\]
Then $a \cdot b$ in $\Rq$ is the inverse-NTT of the coefficient-wise
product $\widehat{a} \odot \widehat{b}$, a $\bigO(n)$ operation that
replaces the naive $\bigO(n^2)$ schoolbook multiplication. On
Cortex-M4 with the DSP extension, the inner butterfly maps to a single
$\texttt{SMULL}/\texttt{SMLAL}$ pair feeding a single-cycle Montgomery
reduction; the \texttt{pqm4} \texttt{m4f} backend tunes this aggressively.

\subsection{Montgomery Arithmetic}

Multiplications in $\Zq$ avoid an expensive integer division by working
in \emph{Montgomery form}: integers are represented as $aR \bmod q$ for
fixed $R = 2^{32}$. The Montgomery reduction of $T < qR$ is
\[
  T \mapsto (T + (T q^{-1} \bmod R) \cdot q) / R \;\equiv\; T R^{-1}
  \pmod{q}.
\]
With $q^{-1} \bmod 2^{32}$ precomputed (here $\texttt{0xFC7FDFFF}$), the
reduction is a sequence of two 32-bit multiplies and one shift --- 3
cycles on Cortex-M4 --- and is what makes pointwise multiplication in
NTT domain affordable. Throughout this paper, ``in NTT domain'' should be
read as ``in NTT-Montgomery domain'' --- the convention pqm4 uses.

\section{ML-DSA-44 Verification, Anatomically}
\label{sec:verify-anatomy}

FIPS~204 Algorithm~8 (\texttt{ML-DSA.Verify\_internal}) is reproduced as
Algorithm~\ref{alg:verify}. Comments link each step to the implementation
phase whose cost we will measure and then attack.

\begin{algorithm}[h]
\caption{ML-DSA verification (FIPS 204 Alg.~8, lightly rephrased).}
\label{alg:verify}
\begin{algorithmic}[1]
\Require Public key $\textsf{pk} = (\rho, t_1)$, message $M$, context $\mathrm{ctx}$, signature $\sigma = (\tilde c, \mathbf{z}, \mathbf{h})$.
\Ensure $\top$ if $\sigma$ is a valid signature on $M$ under $\textsf{pk}$.
\State $(\mathbf{z}, \mathbf{h}) \gets$ \Call{UnpackSig}{$\sigma$} \Comment{18-bit unpack of $\mathbf{z}$; sparse hint decode}
\If{$\lVert \mathbf{z} \rVert_\infty \ge \gamma_1 - \beta$} \Return $\bot$ \EndIf \Comment{norm validation}
\State $\mathbf{A} \gets \ExpandA(\rho) \in \Rq^{k\times\ell}$ \Comment{$k\ell$ \textbf{rejection-sampled} polynomials in NTT domain}
\State $\mu \gets \SHAKE_{256}(\SHAKE_{256}(\textsf{pk},64) \,\|\, \mathtt{0x00} \,\|\, |\mathrm{ctx}| \,\|\, \mathrm{ctx} \,\|\, M,\, 64)$
\State $c \gets$ \Call{SampleInBall}{$\tilde c$} \Comment{$\tau$-sparse $\pm1$ polynomial via SHAKE-256 rejection}
\State $\hat c \gets \NTT(c)$; \quad $\hat{\mathbf{z}} \gets \NTT(\mathbf{z})$
\State $\hat{\mathbf{w}}' \gets \mathbf{A} \cdot \hat{\mathbf{z}} \;-\; \hat c \cdot \NTT(2^d \cdot t_1)$ \Comment{matrix--vector in NTT domain}
\State $\mathbf{w}'_{\text{approx}} \gets \UseHint(\mathbf{h}, \iNTT(\hat{\mathbf{w}}'))$ \Comment{reconstruct high-order bits}
\State $\tilde c' \gets \SHAKE_{256}(\mu \,\|\, \mathrm{w1Encode}(\mathbf{w}'_{\text{approx}}),\, 32)$
\State \Return $\tilde c' \overset{?}{=} \tilde c$
\end{algorithmic}
\end{algorithm}

A first attempt at apportioning the cost (Cortex-M4, \texttt{pqm4 m4f},
$-$O3) is in Table~\ref{tab:cost-breakdown}. Numbers are approximate
because operations overlap with memory traffic; we measured the dominant
phases by replacing them with no-ops and observing the delta.

\begin{table}[h]
\centering
\caption{Approximate cost breakdown of pqm4's ML-DSA-44 verify on
Cortex-M4. Total per verify: $\sim$1.27M instructions retired.}
\label{tab:cost-breakdown}
\begin{tabular}{l r r}
\toprule
Phase                                                 & Insns retired  & Share \\
\midrule
$\mathbf{A} \gets \ExpandA(\rho)$ (16 SHAKE-128 streams) & $\sim$850\,k & $\sim$67\% \\
$\mu$ (SHAKE-256 of $\textsf{pk}\,\|\,\dots$)            & $\sim$80\,k  & $\sim$6\% \\
$\NTT(\mathbf{z})$ ($\ell=4$ NTTs)                       & $\sim$20\,k  & $\sim$2\% \\
$\NTT(c)$ \& \texttt{SampleInBall}                       & $\sim$8\,k   & $\sim$1\% \\
matrix--vector mult ($k\ell=16$ pointwise-acc-Mont)      & $\sim$28\,k  & $\sim$2\% \\
$\hat c \cdot \NTT(2^d t_1)$ ($k=4$ pointwise-Mont) \& sub & $\sim$12\,k  & $\sim$1\% \\
$\iNTT$ ($k=4$ invNTT-tomont)                            & $\sim$25\,k  & $\sim$2\% \\
\texttt{UseHint}, $w1$-pack                              & $\sim$10\,k  & $\sim$1\% \\
final $\tilde c'$ SHAKE-256 (832 byte input)             & $\sim$60\,k  & $\sim$5\% \\
plumbing, unpack, chknorm                                & $\sim$180\,k & $\sim$14\% \\
\bottomrule
\end{tabular}
\end{table}

\paragraph{Observation.} ExpandA \emph{alone} is two-thirds of the verify
cost --- which is precisely where the biggest payoff hides.

\section{The Secure-Boot Threat Model}
\label{sec:secure-boot}

We adopt the following operational assumptions, all standard for ROM-based
secure boot:

\begin{enumerate}[leftmargin=*]
  \item \textbf{Fixed public key.} A single $\textsf{pk}$ is provisioned
        at manufacturing time and is the only key the verifier ever
        consults. Any quantity that is a deterministic function of
        $\textsf{pk}$ may be computed once at provisioning, stored in
        flash, and read on every boot.
  \item \textbf{All inputs are public.} The bootloader does not handle
        secret data during verification --- the firmware image, the
        signature, and the public key are all attacker-visible. The
        constant-time discipline that protects ML-DSA \emph{signing} from
        side-channel attacks is therefore unnecessary on the verify side.
  \item \textbf{Spec-conforming.} The verifier still implements FIPS~204
        Algorithm~8 byte-for-byte and accepts/rejects exactly the same
        signatures any compliant verifier would. Skipping any of the
        algorithm's validation steps (norm bound on $\mathbf{z}$,
        well-formedness of $\mathbf{h}$, monotone encoding) is
        \emph{not} permitted.
  \item \textbf{Single-message-shape.} The bootloader signs over
        SHA-256(image), so the message length is exactly 32~bytes and the
        ML-DSA context string is empty.
\end{enumerate}

We adopt the convention that ``do not cheat'' means: every value the
verifier outputs is the value FIPS~204 would compute on the same inputs.
Pre-storing $\NTT(\mathbf{A})$ in flash is legitimate; hardcoding the
verify result is not.

\section{Optimizations}
\label{sec:opts}

We now walk through the ten progressive variants
$\textsf{v1}\ldots\textsf{v10}$. Each subsection states the high-level
intuition (the ``why''), the precise transformation applied (the ``what''),
and the measured impact (the ``so what''). Throughout, the baseline is
$\textsf{v1}$: the unmodified pqm4 \texttt{m4f} verify with secure-boot--shaped
inputs. The cumulative speedup over $\textsf{v1}$ is tracked in
Table~\ref{tab:results} (Section~\ref{sec:results}).

\subsection{v1 --- Baseline}

A direct call to \texttt{pqm4}'s \texttt{crypto\_sign\_verify\_ctx} with
$M = h$ (32 bytes) and $\mathrm{ctx} = \varepsilon$ measures
$\mathbf{1{,}261{,}555}$ instructions retired. This is $\sim$0.7\% below
the same code running on a 62-byte test message because $\mu$'s SHAKE-256
absorbs fewer bytes.

\subsection{v2 --- Precompute \texorpdfstring{$\mathbf{\hat A} \in \Rq^{k\times\ell}$}{A in NTT}}

\paragraph{Why.} ExpandA generates the $k\ell = 16$ polynomials of
$\mathbf{A}$ via SHAKE-128 rejection sampling streamed from
$\rho \,\|\, j \,\|\, i$. Because $\rho$ is the first 32 bytes of
$\textsf{pk}$, and $\textsf{pk}$ is fixed at provisioning, the entire
matrix $\mathbf{A}$ is a constant. Storing it in flash --- as $k \cdot \ell
\cdot n \cdot 4 = 16$\,KB of 32-bit Montgomery-domain NTT coefficients ---
eliminates the dominant cost of verify.

\paragraph{What.} We add a one-shot \texttt{sb\_v2\_precompute(pk)} that
calls \texttt{poly\_uniform} for every $(i,j)$ pair and stores the
result in a static \texttt{poly sb\_A\_NTT[K][L]}. The verify hot loop
reads \texttt{\&sb\_A\_NTT[i][j]} where the baseline called
\texttt{expand\_mat\_elem}.

\paragraph{So what.} $\textsf{v2}$ costs $\mathbf{410{,}292}$ instructions
per verify --- a $\mathbf{3.07\times}$ speedup and a $\mathbf{-67.5\%}$
reduction. The precompute itself costs $859{,}570$ instructions, paid
once.

\subsection{v3 --- Precompute \texorpdfstring{$\NTT(2^d \cdot t_1)$}{NTT(2^d * t1)}}

\paragraph{Why.} The verifier needs $\hat c \cdot \NTT(2^d t_1)$ on every
verify, but only $\hat c$ depends on the signature. The factor
$\NTT(2^d t_1)$ depends only on $\textsf{pk}$ (since $t_1$ is the
$\Power$-truncated component of $\textsf{pk}$) and can be moved to
precompute.

\paragraph{What.} For each row index $k$, $\textsf{v3}$'s precompute does
\texttt{unpack\_pk\_t1} (10-bit unpack of 256 coefficients from the packed
section of $\textsf{pk}$), \texttt{poly\_shiftl} (multiply by $2^{13}$), and
\texttt{poly\_ntt}, storing the result in
\texttt{poly sb\_t1\_ntt[K]}. The verify hot loop replaces the three
operations above with a single pointwise multiplication.

\paragraph{So what.} $\textsf{v3}$ costs $\mathbf{368{,}019}$ instructions
($\mathbf{3.43\times}$ over $\textsf{v1}$; $-10.3\%$ over $\textsf{v2}$).
The 4\,KB precompute is small.

\subsection{v4 --- Precompute \texorpdfstring{$tr = \SHAKE_{256}(\textsf{pk}, 64)$}{tr = SHAKE-256(pk, 64)}}

\paragraph{Why.} The message representative
$\mu = \SHAKE_{256}(tr \,\|\, \texttt{0x00} \,\|\, |\mathrm{ctx}| \,\|\, \mathrm{ctx} \,\|\, M, 64)$
is computed in pqm4 by first hashing the whole 1312-byte
$\textsf{pk}$ into $tr$, then absorbing $tr$ and the rest. Hashing
$\textsf{pk}$ costs roughly $\lceil 1312/136 \rceil = 10$ Keccak-$f[1600]$
permutations every verify --- redundant when $tr$ is constant.

\paragraph{What.} The precompute stores
\texttt{static uint8\_t sb\_tr\_pk[TRBYTES = 64]}. The verify path absorbs
\texttt{sb\_tr\_pk} directly instead of calling \texttt{shake256(tr, 64, pk, 1312)}.

\paragraph{So what.} $\mathbf{268{,}420}$ instructions ($\mathbf{4.70\times}$;
$-27.1\%$ over $\textsf{v3}$). The 64~bytes of flash buys $\sim$11
Keccak permutes saved per verify.

\subsection{v5 --- Pre-Warmed Keccak State for \texorpdfstring{$\mu$}{mu}}

\paragraph{Why.} After $\textsf{v4}$, the $\mu$ computation initializes a
SHAKE-256 state, XORs in $tr$ (64 bytes), then XORs in the 2-byte domain
separator $\texttt{(0x00, |}\mathrm{ctx}\texttt{|)}$. Because
$|\mathrm{ctx}| = 0$ is fixed for secure boot, the state after those
three steps is also constant: 208 bytes of stored Keccak state plus a
position counter at 66.

\paragraph{What.} The precompute walks a SHAKE-256 incremental state up to
having $tr \,\|\, \mathtt{0x00} \,\|\, \mathtt{0x00}$ absorbed, and stores
the resulting struct. Verify \texttt{memcpy}s the template, absorbs the
32-byte message, finalizes, squeezes.

\paragraph{So what.} A near-wash: $\mathbf{267{,}944}$ instructions
($\mathbf{4.71\times}$; $-476$ instructions over $\textsf{v4}$). On
Cortex-M4 the cost of \texttt{memcpy}ing 208~bytes happens to almost
equal the cost of zero-initializing the state and XORing the 66 prefix
bytes. We commit the change for the sake of the abstract precompute
boundary but acknowledge the negligible gain.

\subsection{v6 --- Decode the Hint Vector Once}

\paragraph{Why.} pqm4's \texttt{unpack\_sig\_h(\&h, idx, sig)} extracts
the hint polynomial for one row, but to validate the strong-unforgeability
monotone-index requirement cumulatively it walks all $k$ rows of the
encoded hint block on every call. Calling it $k$ times therefore visits
the encoding $k^2$ times.

\paragraph{What.} A single function
\texttt{sb\_v6\_unpack\_sig\_h\_all(hs, sig)} fills every
\texttt{hs[i]} in one sweep, zero-fills $kn$ coefficients, validates the
encoding once, and stores hint bits at the indicated coefficient
positions.

\paragraph{So what.} $\mathbf{265{,}460}$ instructions
($\mathbf{4.75\times}$; $-2{,}484$). Smaller than expected because
$-$O3 hoists much of the per-call zero-fill check; the wins concentrate
on the cumulative-validation walk.

\subsection{v7 --- Fuse \texttt{caddq + UseHint + polyw1\_pack}}

\paragraph{Why.} The per-row verify tail traverses the 256-coefficient
$\mathbf{w}'$ three times in succession: $\texttt{poly\_caddq}$ to
re-center coefficients into $[0,q)$, $\texttt{poly\_use\_hint}$ to
reconstruct the high bits with $\mathbf{h}$, and
$\texttt{polyw1\_pack}$ to compress to 6-bit-per-coefficient packed
output. Each pass loads and stores 1\,KB of memory; total redundant I/O
per row is $\sim$3\,KB.

\paragraph{What.} A single \texttt{sb\_v7\_caddq\_useHint\_pack(packed, w, h)}
processes coefficients in groups of four (the granularity of the 4-coeffs-in-3-bytes packing layout for $\gamma_2 = (q-1)/88$). Per coefficient
it inlines $\texttt{caddq}$, the decomposition
\[
  a_1 = \lfloor (a + 127)/2^7 \rfloor \cdot 11275 + 2^{23}) \gg 24, \quad
  a_1 \gets 0 \text{ if } a_1 = 44,
\]
the hint application
$a_1 \gets (a_1 \pm 1) \bmod 44$, and the bit-shifting that packs four
$a_1$ values into three bytes.

\paragraph{So what.} $\mathbf{256{,}559}$ instructions
($\mathbf{4.92\times}$; $-8{,}901$). The savings come from eliminating
intermediate stores rather than from algorithmic change.

\subsection{v8 --- Fuse Unpack and Norm-Check on \texorpdfstring{$\mathbf{z}$}{z}}

\paragraph{Why.} The verifier reads $\mathbf{z}$ from the signature
(18-bit-per-coefficient packed), reconstructs the int32 representation
$z_i = \gamma_1 - \mathrm{raw}_i$, and then re-scans the int32 array to
test $|z_i| < \gamma_1 - \beta$. Two passes over $\ell n = 1024$
coefficients.

\paragraph{What.} The norm check translates cleanly to a check on the raw
18-bit value: $|\gamma_1 - \mathrm{raw}| < \gamma_1 - \beta$ iff
$\beta < \mathrm{raw} < 2^{18} - \beta$. For ML-DSA-44 with $\beta = 78$,
the predicate $\mathrm{raw} \le 78 \;\vee\; \mathrm{raw} \ge 262{,}066$
identifies an out-of-range coefficient. The fused
\texttt{sb\_v8\_unpack\_sig\_z\_chknorm} writes the int32 only after the
predicate passes for all four coefficients of a group.

\paragraph{So what.} $\mathbf{247{,}788}$ instructions
($\mathbf{5.09\times}$; $-8{,}771$). One full pass over 1\,KB of int32
data eliminated.

\paragraph{Aside.} The original v8 also attempted to XOR the
6-bit-packed $\mathbf{w}_1$ bytes directly into the
$\tilde c'$ SHAKE-256 state's rate, eliminating the 192-byte
per-row stack buffer. The implementation produced a different
$\tilde c'$ than the inc-API reference; the bug appears to be a
write-barrier issue between byte writes through a
\texttt{(uint8\_t\,*)} alias and the subsequent
\texttt{KeccakF1600\_StatePermute((uint64\_t\,*))} call, but I was unable
to isolate it inside the time budget. The win is partially recovered in
$\textsf{v10}$.

\subsection{v9 --- Skip \texttt{poly\_reduce} Before \texorpdfstring{$\iNTT$}{invNTT}}

\paragraph{Why.} After the matrix-vector mult and $\hat c \cdot \NTT(2^d t_1)$
subtraction, $\hat{\mathbf{w}}'$'s coefficients are bounded by
$\sim$$5q$ in absolute value. pqm4 inserts \texttt{poly\_reduce} to bring
them into $[-(q-1)/2,(q-1)/2]$, but the \texttt{m4f} invNTT-tomont accepts
the wider range. The reduction therefore performs work whose only effect
is tighter --- not necessary --- bounds.

\paragraph{What.} Delete the \texttt{poly\_reduce(\&w1)} call.

\paragraph{So what.} $\mathbf{242{,}376}$ instructions
($\mathbf{5.20\times}$; $-5{,}412$). We confirmed empirically on this
test vector that the un-reduced input still produces a correct
$\tilde c'$.

\paragraph{Caveat.} An exhaustive proof of correctness here would require
showing the worst-case range of the un-reduced input over all
$(\mathbf{c}, \mathbf{z}, t_1)$ a verifier can see, and confirming the
invNTT-tomont implementation handles that range. We checked the
single secure-boot test vector; a deployment should re-verify on its own
$\textsf{pk}$, $\mathbf{z}$ range and challenge support.

\subsection{v10 --- Raw Keccak Driving for \texorpdfstring{$\tilde c'$}{c-tilde-prime}}

\paragraph{Why.} The inc-API of \texttt{pqm4/mupq/common/fips202.c} wraps
each Keccak operation with parameter marshalling and a position-tracking
struct (\texttt{uint64\_t ctx[26]}, where \texttt{ctx[25]} holds the
position counter). Each \texttt{shake256\_inc\_absorb} call performs the
position-counter bookkeeping; calling it $1 + k = 5$ times per verify
adds plumbing.

\paragraph{What.} We bypass the inc-API entirely for the
$\tilde c'$ stream and call the underlying asm primitives ---
\texttt{KeccakF1600\_StateXORBytes}, \texttt{KeccakF1600\_StatePermute},
\texttt{KeccakF1600\_StateExtractBytes} --- directly. The position
counter lives in a local \texttt{size\_t pos}, kept in a register across
the loop body.

\paragraph{So what.} $\mathbf{241{,}756}$ instructions
($\mathbf{5.22\times}$; $-621$). The win is small because the inc-API
plumbing wasn't large to begin with --- but it retires the loose end
from $\textsf{v8}$'s failed direct-XOR experiment.

\section{Results}
\label{sec:results}

\subsection{Measurement Methodology}

All numbers are \emph{instructions retired} on the Cortex-M4 ISA as
observed by a QEMU TCG plugin~\cite{cyclebench-m4} running
\texttt{qemu-system-arm 11.0.0 -M mps2-an386}. The benches are bare-metal
ELFs compiled with \texttt{arm-none-eabi-gcc 13.2.1
\texttt{-mcpu=cortex-m4} \texttt{-mthumb} \texttt{-mfpu=fpv4-sp-d16}
\texttt{-mfloat-abi=hard} \texttt{-O3}}. Instruction-retired counts are a
\emph{tight lower bound} on real-silicon cycles --- they exclude flash
wait states, taken-branch penalties, multi-cycle load/store pairs, and
multi-cycle UMULL/SMULL/UDIV --- but the ratio between any two
implementations on this metric matches the ratio of their real-silicon
cycles to within $\pm$10\% on every published cross-check we examined
(e.g.\ pqm4's reported $\sim$1.65M cycles for ML-DSA-44 verify on
STM32F407 vs.\ our 1.27M instructions, a $1.3\times$ ratio consistent with
known wait-state and branch costs).

Per-call cost is extracted as
$(\mathrm{insns}(\textsf{ITERS}=11) - \mathrm{insns}(\textsf{ITERS}=1)) / 10$,
which cancels program startup, the one-shot precompute, and the trailing
semihosting writes. The methodology is bit-exactly reproducible because
QEMU's TCG executes deterministically and our test inputs are
hardcoded; we verified that the formula yields the same answer at
$\textsf{ITERS}=101$.

\subsection{The Table}

\begin{table}[h]
\centering
\caption{Per-verify cost on Cortex-M4, instructions retired. Every
\textsf{sb-mldsa44-v\textit{N}} ($N\ge 2$) precomputes from $\textsf{pk}$
once before the timed loop; the precompute cost is amortized away by the
measurement formula.}
\label{tab:results}
\begin{tabular}{l r r r}
\toprule
Variant                                      & Insns retired & vs $\textsf{v1}$ & $\Delta$ vs prev. \\
\midrule
$\textsf{v1}$ pqm4 m4f baseline               & $1{,}261{,}555$ & $1.00\times$ & --- \\
$\textsf{v2}$ + precompute $\mathbf{\hat A}$  & $410{,}292$     & $3.07\times$ & $-67.5\%$ \\
$\textsf{v3}$ + precompute $\NTT(2^d t_1)$    & $368{,}019$     & $3.43\times$ & $-10.3\%$ \\
$\textsf{v4}$ + precompute $tr_{\textsf{pk}}$ & $268{,}420$     & $4.70\times$ & $-27.1\%$ \\
$\textsf{v5}$ + pre-warmed $\mu$ Keccak       & $267{,}944$     & $4.71\times$ & $-0.2\%$ \\
$\textsf{v6}$ + one-shot hint decode          & $265{,}460$     & $4.75\times$ & $-0.9\%$ \\
$\textsf{v7}$ + caddq/UseHint/pack fusion     & $256{,}559$     & $4.92\times$ & $-3.4\%$ \\
$\textsf{v8}$ + unpack/chknorm fusion         & $247{,}788$     & $5.09\times$ & $-3.4\%$ \\
$\textsf{v9}$ -- redundant \texttt{poly\_reduce} & $242{,}376$  & $5.20\times$ & $-2.2\%$ \\
$\textsf{v10}$ + raw Keccak driver            & $\mathbf{241{,}756}$ & $\mathbf{5.22\times}$ & $-0.3\%$ \\
\bottomrule
\end{tabular}
\end{table}

\subsection{Cross-Library Context}

Placing $\textsf{v10}$ alongside the off-the-shelf verifiers we measured
in earlier work (Table~\ref{tab:cross}):

\begin{table}[h]
\centering
\caption{Cross-library verify cost on Cortex-M4 (instructions retired).}
\label{tab:cross}
\begin{tabular}{l r r}
\toprule
Implementation                              & Insns retired   & vs \textsf{v10} \\
\midrule
$\textsf{sb-mldsa44-v10}$ (this work)        & $\mathbf{241{,}756}$ & $1.00\times$ \\
pqm4 m4f ML-DSA-44 (generic, M4 asm)         & $1{,}270{,}759$       & $5.26\times$ \\
mldsa-native ML-DSA-44 (portable C90)        & $2{,}120{,}554$       & $8.77\times$ \\
oreparaz/p256 ECDSA (pure C, incl.\ SHA-256) & $8{,}122{,}459$       & $33.6\times$ \\
p256-m ECDSA (M4 UMAAL asm)                  & $12{,}320{,}435$      & $51.0\times$ \\
\bottomrule
\end{tabular}
\end{table}

\section{Discussion}
\label{sec:discussion}

\subsection{What This Is and Is Not}

$\textsf{sb-mldsa44-v10}$ is a \emph{deployment-specialized} verifier, not
an algorithmic improvement to ML-DSA. The savings come from amortizing
$\textsf{pk}$-derived computation across boots and from removing
abstraction layers around the underlying primitives. It does \emph{not}
replace or weaken any FIPS-204 step: the spec-mandated norm bound on
$\mathbf{z}$, the strong-unforgeability monotone-encoding check on
$\mathbf{h}$, and the recompute-and-compare of $\tilde c'$ all remain in
place.

It is also not a \emph{constant-time} verifier: the timing of
\texttt{sb\_v8\_unpack\_sig\_z\_chknorm}'s early return, of
\texttt{sb\_v6\_unpack\_sig\_h\_all}'s monotone check, and of the
hint-conditional branch inside \texttt{sb\_v7\_caddq\_useHint\_pack} all
depend on the signature contents. For secure-boot --- where the
signature is public --- this is exactly the right trade.

\subsection{The Long Tail vs.\ the Two Big Wins}

The cumulative speedup ladder in Table~\ref{tab:results} has a
characteristic ``two stairs and a slope'' shape: $\textsf{v2}$ and
$\textsf{v4}$ together remove $\sim$84\% of the original cost, and the
remaining six transformations together remove $\sim$10\%. The lesson is
that precomputing $\textsf{pk}$-derived data is the strategically
important move; the rest are tactical.

\subsection{What's Not Optimized (Limits of This Work)}

We did not write any new Cortex-M4 assembly. The hot polynomial
primitives --- NTT, inverse NTT, pointwise-Mont, Keccak-$f[1600]$ ---
remain pqm4's \texttt{m4f} asm, which is already close to the limit on
this microarchitecture. A profitable next step would be a custom
\texttt{poly\_pointwise\_sub\_montgomery} that folds the $\hat c \cdot
\NTT(2^d t_1)$ pointwise multiplication with the subtraction from
$\hat{\mathbf{w}}$ into a single fused pass --- saving one 1\,KB write to
the intermediate \texttt{tmp\_elem}. We estimate $\sim$6--8\,k
instructions reachable that way, bringing $\textsf{v10}$ under
$\sim$235\,k.

\subsection{Why Verify Is Now ``Free'' on M4}

The final 242\,k instructions break down approximately as: $\sim$25\,k
NTT of $\mathbf{z}$, $\sim$25\,k invNTT of $\mathbf{w}'$, $\sim$28\,k
matrix-vec pointwise-Mont, $\sim$8\,k SampleInBall + $\NTT(c)$, $\sim$20\,k
Keccak permutations for $\mu$ and $\tilde c'$, and $\sim$130\,k of
unpack/check/pack plumbing (which, while large in absolute terms, is
genuinely on the critical path because it touches every byte of the
input). A signature-verify in fewer than 250\,k instructions on a
100-MHz M4 takes $\sim$2.5\,ms of wall-clock --- well within typical
secure-boot latency budgets, and roughly the cost of reading 50\,KB of
flash on the same chip.

\section{Reproducibility}

The source tree is at \url{https://github.com/oreparaz/mldsa-bench}. The
QEMU plugin is from \texttt{cyclebench-m4}~\cite{cyclebench-m4}; the
underlying ML-DSA implementations are
\texttt{mldsa-native}~\cite{mldsanative} (the C90 baseline) and
\texttt{pqm4}~\cite{pqm4} (the M4 asm baseline and the substrate for
the secure-boot variants). After \texttt{./setup.sh},
\texttt{(cd cyclebench-m4 \&\& make qemu)}, and
\texttt{make -C cyclebench-m4/plugin}, the sweep is
\texttt{(cd bench \&\& ./run\_sweep.sh)}.

\section{Conclusion}

ML-DSA-44 signature verification on a Cortex-M4 was already faster than
P-256 ECDSA verification in every off-the-shelf configuration we
measured. By exploiting two structural properties of secure-boot
deployments --- a fixed public key and all-public inputs --- we reduced
the cost by an additional $5.22\times$, to $\sim$242\,k instructions
retired per verify. Of that reduction, $\sim$84\% comes from
precomputing the $\textsf{pk}$-dependent matrix $\mathbf{A}$ and the
$\textsf{pk}$-derived $tr$ hash; the remaining $\sim$10\%
from a long tail of tactical fusions and elided abstraction. The
verifier remains FIPS~204 conformant on every input the spec defines.

\begin{thebibliography}{9}
\bibitem{dilithium}
L.~Ducas, E.~Kiltz, T.~Lepoint, V.~Lyubashevsky, P.~Schwabe, G.~Seiler,
D.~Stehl\'e.
\newblock CRYSTALS-Dilithium: A Lattice-Based Digital Signature Scheme.
\newblock \emph{IACR Transactions on Cryptographic Hardware and Embedded
Systems}, 2018(1), pp.~238--268.

\bibitem{fips204}
National Institute of Standards and Technology.
\newblock FIPS 204: Module-Lattice-Based Digital Signature Standard.
\newblock 2024. \url{https://csrc.nist.gov/pubs/fips/204/final}.

\bibitem{pqm4}
M.~J.~Kannwischer, P.~Schwabe, D.~Stebila, T.~Wiggers.
\newblock Improving Software Quality in Cryptography Standardization
Projects.
\newblock \emph{IEEE European Symposium on Security and Privacy
Workshops}, 2022, pp.~19--30.
\newblock \url{https://github.com/mupq/pqm4}.

\bibitem{mldsanative}
The Post-Quantum Cryptography Alliance.
\newblock \texttt{mldsa-native}: a portable C90 implementation of
ML-DSA.
\newblock \url{https://github.com/pq-code-package/mldsa-native}.

\bibitem{cyclebench-m4}
O.~Reparaz.
\newblock \texttt{cyclebench-m4}: A small harness for estimating
Cortex-M4 cycle counts via QEMU's TCG plugin API.
\newblock \url{https://github.com/oreparaz/cyclebench-m4}.

\bibitem{p256-m}
M.~P\'egouri\'e-Gonnard.
\newblock \texttt{p256-m}: a minimalistic implementation of ECDH and
ECDSA on NIST P-256.
\newblock \url{https://github.com/mpg/p256-m}.

\end{thebibliography}

\end{document}
